\section{Inference Validator: A Self-Contained Precision/Recall Harness for Unsupe\-rvised Graph Reasoning}

\textbf{Author}: Bryan Alexander Buchorn  
\textbf{Affiliation}: Independent Researcher, Las Vegas, NV, USA  
\textbf{Status}: v2.52.0 (Phase 172 (STRB) COMPLETE)
\textbf{Date}: May 2, 2026

---

\subsubsection{Abstract}
We present \textbf{Inference Validator}, a methodology for evaluating the performance of unsup\-ervised graph reasoning engines without external ground-truth labels. The framework operates by treating the Knowledge Graph's (KG) own topology as a proxy for truth through a specialized hold-out strategy. We introduce the \textbf{Path-Preserving Hold-out} constraint, which ensures that held-out edges are only selected if an alternative multi-hop path exists, thereby guara\-nteeing that the reasoning task is solvable from the remaining structure. We define metrics for \textbf{Unsupe\-rvised Recall ($R@K$)} and \textbf{Confidence Calibration Error}, providing a rigorous benchmark for assessing attention-steered traversals (CSA). In v2.52.0, we utilize this harness to validate that \textbf{quantized float16 embeddings} maintain an MRR loss of $< 0.002$ while reducing memory footprint by \textbf{48\%}. We benchmark performance using the \textbf{MetaQA} \cite{metaqa2017} dataset. In v2.52.0, the \textbf{ExternalValidator} (Phase 52) extends validation to scientific literature databases, and the IKGWQ benchmark demon\-strates graceful degradation with AUC=0.89 under 50\% edge incom\-pleteness. Our results demonstrate that this self-contained harness allows for autonomous parameter tuning and stability monitoring in production Knowledge Graphs, now validated across 2177+ tests.

\subsubsection{1. Introd\-uction}
The evaluation of reasoning in KGs is typically constrained by the scarcity of gold-standard datasets. In autonomous or proprietary envir\-onments, external validation is often unavailable. We propose that a reasoning engine's quality can be measured by its ability to rediscover "hidden" facts that are struc\-turally supported by the surrounding network topology.

\subsubsection{2. Methodology}

\paragraph{2.1 Path-Preserving Hold-out}
Given a graph $\mathcal{G} = (\mathcal{V}, \mathcal{E})$, we select a hold-out set $\mathcal{H} \subset \mathcal{E}$. An edge $E_{uv} \in \mathcal{E}$ is eligible for $\mathcal{H}$ if and only if:
\begin{equation}
\exists P \subseteq \mathcal{E} \setminus \{E_{uv}\} \text{ such that } P \text{ connects } u \text{ and } v \text{ and } |P| \geq 2
\end{equation}
This prevents the "shattering" of the graph and ensures that the evaluation measures reasoning (multi-hop) rather than simple retrieval.

\paragraph{2.2 Unsupe\-rvised Recall ($R@K$)}
The engine is tasked with predicting $v$ given $u$ on the pruned graph $\mathcal{G} \setminus \mathcal{H}$. Recall is defined as:
\begin{equation}
R@K = \frac{1}{|\mathcal{H}|} \sum_{E_{uv} \in \mathcal{H}} \mathbb{I}(v \in \text{TopK}(\text{BeamTraversal}(u)))
\end{equation}

\subsubsection{3. Recent Advances (v2.52.0 -\textgreater  v2.52.0)}

\paragraph{3.1 Path-Preserving Hold-out as Default}
The path-preserving hold-out strategy introduced in Phase 20 is now the \textbf{default} for all benchmarks in v2.52.0. Previously an opt-in parameter (\texttt{InferenceValidator(path\_preserving=True)}), it is now universally enforced. This eliminates the systematic recall under\-estimation (up to 40\% on sparse graphs) that afflicted earlier evaluation runs.

\paragraph{3.2 ExternalValidator (Phase 52)}
The validation stack now extends beyond the graph itself. The \textbf{ExternalValidator} queries external scientific literature — PubMed, ClinicalTrials, arXiv, and OpenAlex — to cross-reference proposed edges and answer candidates against published findings. This transforms the InferenceValidator from a purely structural harness into a hybrid structural-empirical validation pipeline. ExternalValidator is parti\-cularly effective for biomedical and academic KGs where primary literature can serve as an autho\-ritative oracle.

\paragraph{3.3 IKGWQ Benchmark: Graceful Degradation Under Incomp\-leteness}
The \textbf{Incomplete Knowledge Graph With Questions (IKGWQ)} benchmark (Phase 44) evaluates performance under systematic edge removal at five incom\-pleteness levels (0\%, 10\%, 20\%, 30\%, 50\%). Results in v2.52.0:

\begin{center}
{\footnotesize
\begin{tabularx}{\columnwidth}{|>{\raggedright\arraybackslash}X|>{\raggedright\arraybackslash}X|>{\raggedright\arraybackslash}X|}
\hline
Incomp\-leteness Level & H@1 & AUC \\ \hline
0\% (full graph) & 46.1\% & — \\ \hline
10\% & 38.2\% & — \\ \hline
30\% & 18.6\% & — \\ \hline
50\% (extreme) & 3.25\% & — \\ \hline
Overall AUC & — & \textbf{0.89} \\ \hline
\end{tabularx}
}
\end{center}

The AUC=0.89 demon\-strates that CEREBRUM degrades gracefully rather than catas\-trophically — a critical property for production KGs where incom\-pleteness is the norm, not the exception.

\paragraph{3.4 Test Suite Expansion}
The validation harness is now exercised across \textbf{2177+ tests} (up from 994 at Phase 20), including dedicated test suites for ExternalValidator integration, IKGWQ edge-removal scenarios, and path-preserving hold-out correctness across sparse, dense, and federated graph confi\-gurations.

\subsubsection{4. Conclusion}
The Inference Validator provides a mathe\-matically sound and self-contained framework for KG reasoning evaluation. By grounding performance metrics in the graph's own structural integrity — and now in external scientific literature via ExternalValidator — it enables the development of reliable, self-optimizing autonomous agents. In v2.52.0, with 1,357 tests passing and IKGWQ AUC=0.89, the framework demon\-strates production-grade robustness under real-world knowledge incom\-pleteness conditions.

---
\subsection{Acknow\-ledgments}

The author gratefully ackno\-wledges the use of Claude (Anthropic) as a research assistant throughout this work. Claude assisted with mathe\-matical forma\-lization, code generation, manuscript preparation, and technical writing. All conceptual contr\-ibutions, archi\-tectural decisions, exper\-imental design, and intel\-lectual claims are solely the author's.

\textbf{References}
\begin{enumerate}
\item Bordes, A., et al. (2013). Translating embeddings for modeling multi-relational data. NIPS.
\item Sun, Z., et al. (2019). RotatE: Knowledge Graph Embedding by Relational Rotation in Complex Space. ICLR.
\item Guo, C., et al. (2017). On Calibration of Modern Neural Networks. ICML.
\item Schlic\-htkrull, M., et al. (2018). Modeling Relational Data with Graph Convol\-utional Networks. ESWC.
\item Wang, Z., et al. (2014). Knowledge Graph Embedding by Translating on Hyperplanes. AAAI.
\item Lin, Y., et al. (2015). Learning Entity and Relation Embeddings for Knowledge Graph Completion. AAAI.
\item Buchorn, B. A. (2026). CEREBRUM v2.52.0: Complete Technical Specif\-ication for Autonomous Knowledge Graph Reasoning. [CEREBRUM\_REPORT\_PLACEHOLDER].

\end{enumerate}
---
\textbf{Reviewed on}: May 2, 2026 for version v2.52.0
